\documentclass[11pt,reqno]{amsart}
\input{auxiliaryFile.tex}
%JournalMIM_Template_v2.3.3 / 19.07.2024

%%insert new packages/commands here if needed
%\usepackage{}
%\usepackage{}
%\newcommand{\grad}{\mathop{\mathrm{grad}}}


%%Write the Orcid IDs below in order of authors, such as OrcidA for the 1st author, OrcidB for the 2nd author, and so on.
\newcommand{\orcidauthorA}{0000-0001-0000-1111}
\newcommand{\orcidauthorB}{0000-0002-0000-2222}
%\newcommand{\orcidauthorC}{0000-0003-0000-XXXX}

\begin{document}
\scalebox{0.3}{\includegraphics[]{logo.png}}

\vspace{-2.6cm}
\hfill{\bf International Journal of Maps in Mathematics}

\hfill{\it  Volume \textbf{X}, Issue \textbf{X}, 202., Pages:\textbf{XX-XX}}

\hfill{E-ISSN: 2636-7467}

\hfill{\color{blue}{\bf www.journalmim.com}}

\setcounter{page}{1}
\linenumbers

\vspace{1cm}
\title[Int. J. Maps Math. (202X) XX(X):xx--xx / Short Title]{Full Title}
\author[F. Author]{First Author \orcidA{} $^*$} 
\address{First Author's Address}
\author[S. Author]{Second Author \orcidB{}}
\address{Second  Author's Address}
%\author[T. Author]{Third Author \orcidC{}} %\address{Third Author's Address}
%add new rows if needed

%add new rows if it is needed
\thanks{
{\it Received:  \hspace{2.5cm}Revised:  \hspace{2.5cm} Accepted:}\\
$^*$ Corresponding author\\
Name of the 1st Author $\diamond$ Email of 1st Author $\diamond$ https://orcid.org/\orcidauthorA\\
Name of the 2nd Author $\diamond$ Email of 2nd Author $\diamond$ https://orcid.org/\orcidauthorB\\
%Name of the 3rd Author $\diamond$ Email of 3rd Author $\diamond$ https://orcid.org/\orcidauthorC\\
}

\begin{abstract}
	\hrule width15cm\bigskip 
The manuscripts will include the full address (es) of the author (s), with E-mail address (es) and ORCID id(s), an abstract not exceeding 300 words, 2010 Mathematics Subject Classification, Key words and phrases. All illustrations, figures, and tables are placed within the text at the appropriate points, rather than at the end.\\  %Please replace this text with your abstract
	\textbf{Keywords}: Keyword1, Keyword2, ... \\
	\textbf{2010 Mathematics Subject Classification}: Primary, Secondary.\\
	\hrule
\end{abstract}
\maketitle

\vspace{-0.5cm}
\section{Introduction}  %% Please avoid complicated
\lipsum[1]%Please replace this command/text with your content

\begin{theorem} The square of any real number is non-negative.
\end{theorem}

\begin{proof}
	Any real number $x$ satisfies $x>0$, $x=0$, or $x<0$.
	If $x=0$, then $x^2=0\ge 0$.  If $x>0$ then as a positive time a
	positive is positive we have $x^2=xx>0$.  If $x<0$ then $-x>0$ and so
	by what we have just done $x^2=(-x)^2>0$.  So in all cases $x^2\ge0$.
\end{proof}

\begin{definition}\label{def1}
	content...
\end{definition}


\begin{example}\label{ex1}
	content...
\end{example}

\section{Preliminaries}

\lipsum[2]%Please replace this command/text with your content

\begin{table}[h]
	\caption{Caption text}\label{tab1}%
	\begin{tabular}{| l | c c c |}
		\hline
		Column 1 & Column 2  & Column 3 & Column 4\\ \hline
		row 1    & data 1   & data 2  & data 3  \\
		row 2    & data 4   & data 5  & data 6  \\
		row 3    & data 7   & data 8  & data 9  \\	\hline
	\end{tabular}
\end{table}

\lipsum[3]%Please replace this command/text with your content

\bigskip

\begin{equation}\label{eq1}
e^{i\pi}+1=0
\end{equation}

\begin{theorem}\label{teo2} Euler's identity (also known as Euler's equation) is the equality
$ e^{i\pi}+1=0 $
	where
	$e$  is Euler's number, the base of natural logarithms,
	$i$  is the imaginary unit, which by definition satisfies $i^{2}=-1$, and
	$ \pi $ is pi, the ratio of the circumference of a circle to its diameter.
\end{theorem}

\begin{proof}
Please write proof of the Theorem \ref{teo2} here \cite{bib9}.
\end{proof}

\begin{corollary}\label{cor1}
	content...
\end{corollary}

\begin{proposition}
	content...
\end{proposition}


\lipsum[6]%Please replace this command/text with your content

The well known Pythagorean theorem \(x^2 + y^2 = z^2\) was 
proved to be invalid for other exponents. 
Meaning the next equation has no integer solutions:
\[ x^n + y^n = z^n \]

\begin{proof}[Proof of Corollary \ref{cor1}]
	Please write proof of the Corollary \ref{cor1} here \cite{gursoy2022construction}.
\end{proof}

\begin{lemma}
	content...
\end{lemma}

\begin{remark}
	content...
\end{remark}


\lipsum[7]%Please replace this command/text with your content

\section{Conclusion}
\lipsum[3]%Please replace this command/text with your content

\begin{equation*}
	x = a_0 + \cfrac{1}{a_1 
		+ \cfrac{1}{a_2 
			+ \cfrac{1}{a_3 + \cfrac{1}{a_4} } } }
\end{equation*}

\bigskip

\lipsum[2]%Please replace this command/text with your content



%%Acknowledgments

%\textbf{Acknowledgments.} The authors would like to thank the referee for some useful comments and their helpful suggestions that have improved the quality of this paper.



%%References

%If you have a bibliography file with .bib extension, you can create the references section with the following commands \bibliography{} and \bibliographystyle{apa}. 
%\bibliography{bibfile}
%\bibliographystyle{apa} % all references must be in apa format


%If you do not use a bibliography file, create the references section with \begin{thebibliography} as below.
\begin{thebibliography}{99} %% n is number of items, or the largest label
%Please use APA style bibliography

\bibitem{Dasilva} Cannas da Silva, A. (2008). Lectures on symplectic geometry. Vol. 1764. Lecture Notes in Mathematics. Springer-Verlag, Berlin.

\bibitem{Chenbooklast} Chen, B. Y. (2017). Differential geometry of warped product manifolds and submanifolds. Singapore: World Scientific.

\bibitem{Datta-Islam} Datta, M., \& Islam, M. R. (2009). Submersions on open symplectic manifolds. Topology and its Applications, 156(10), 1801-1806.

\bibitem{FP} Falcitelli, M., Ianus, S., \& Pastore, A. M. (2004). Riemannian submersions and related topics, World Sci. Publishing, River Edge, NJ.

\bibitem{Gray} Gray, A. (1967). Pseudo-Riemannian almost product manifolds and submersions. J. Math. Mech., 16, 715-738.

\bibitem{gursoy2022optimization} Gürsoy, A. (2022). Optimization of product switching processes in assembly lines. Arabian Journal for Science and Engineering, 47(8), 10085-10100.

\bibitem{gursoy2022construction} Gürsoy, A. (2022). Construction of networks by associating with submanifolds of almost Hermitian manifolds. Fundamental Journal of Mathematics and Applications, 5(1), 21-31.

\bibitem{Hogan} Hogan, P. A. (1984). Kaluza–Klein theory derived from a Riemannian submersion. Journal of mathematical physics, 25(7), 2301-2305.

\bibitem{O'Neill} O'Neill, B. (1966). The fundamental equations of a submersion. The Michigan Mathematical Journal, 13(4), 459-469.

\bibitem{Sahinbook} Sahin, B. (2017). Riemannian submersions, Riemannian maps in Hermitian geometry, and their applications. Elsevier.

\bibitem{bib9} Sahin, B. (2020). Symplectosubmersions. International Journal of Maps in Mathematics-IJMM, 3(1), 3-9.

\bibitem{Watson}Watson, B. (1976). Almost hermitian submersions. Journal of Differential Geometry, 11(1), 147-165.

\bibitem{Yano-Kon} Yano, K., \& Kon, M. Structures on manifolds, World Scientific (1984). Department of Mathematics, University of California at Riverside, Riverside CA, 92521.


\end{thebibliography}



\end{document}



